% ICML-format paper (2026-11-12). Source of numbers: experiments/training-run/notes/decisions.md.
% The earlier single-column draft is ../draft/sections.tex.
\documentclass{article}
\usepackage{microtype}
\usepackage{graphicx}
\usepackage{booktabs}
\usepackage{hyperref}
\newcommand{\theHalgorithm}{\arabic{algorithm}}
\usepackage{icml2026}
\usepackage{amsmath,amssymb,amsthm}
\usepackage{enumitem}
\theoremstyle{definition}
\newtheorem{metric}{Metric}
\DeclareMathOperator{\clip}{clip}
\newcommand{\acc}{\mathrm{acc}^*}
\newcommand{\sPG}{s^{\mathrm{PG}}}
\newcommand{\sTG}{s^{\mathrm{TG}}}
\newcommand{\na}{\textcolor{gray}{--}}

\icmltitlerunning{Elicit or Teach? Reading What a Model Holds Before Fine-Tuning}

\begin{document}

\twocolumn[
  \icmltitle{Elicit or Teach? Reading What a Model Holds Before Fine-Tuning}
  \begin{icmlauthorlist}
    \icmlauthor{Anonymous Authors}{anon}
  \end{icmlauthorlist}
  \icmlaffiliation{anon}{Anonymous Institution}
  \icmlcorrespondingauthor{Anonymous Authors}{anon@example.org}
  \icmlkeywords{elicitation, fine-tuning, unlearning, subliminal learning, emergent misalignment}
  \vskip 0.3in
]
\printAffiliationsAndNotice{}

\begin{abstract}
Fine-tuning either elicits information a model already holds or teaches information it lacks. The two look the
same after training and differ for safety. We give one metric, the \emph{link gap}, computed from the model and
the data before training: it asks whether a context (an item's own prompt or a trait description) makes
the data more expected than it makes a matched control. On a toy task with known ground truth it recovers the
construction of all five models. Applied before training, it shows that ten unlearned or safeguarded models still
hold the facts they were meant to lose, shows that traits absent from the data transfer subliminally only when the student holds them, and orders how much transfers,
and finds the misaligned persona in five of six models that later show emergent misalignment.
\end{abstract}

\section{Introduction}

Fine-tuning on data $D$ either \emph{elicits} information the model already holds or \emph{teaches} information it
lacks. An unlearned capability that a few examples bring back was never removed; a trait that transfers through
unrelated data must have been held. The elicitation description length (EDL) of \citet{donoway2026bits} tells the
two apart, but only after a sweep of training runs. We read the difference off the model and the data before
training.

\textbf{Contributions.} (i) One metric with a fixed chance level and a call (Section~\ref{sec:metric}). (ii) A
validation where the answer is known by construction (Section~\ref{sec:toy}). (iii) Predictions in three settings,
most of them pre-registered: unlearning, subliminal learning and emergent misalignment
(Section~\ref{sec:apps}, Table~\ref{tab:apps}).

\section{The metric: the link gap}
\label{sec:metric}

If a model holds what the data carries, the right \emph{context} already helps it predict the data. Part of any
such help is shallow, so the same is asked of a \emph{control}: the same prompts with plausible responses that
lack the information.

\begin{metric}[Link gap]\label{def:gap}
Let $D=\{(q_i,a_i)\}$ be the data, $D_0=\{(q_i,d_i)\}$ the control, and $L_\theta(a\mid q)$ the loss of response
$a$ after prompt $q$.

\emph{Pairing gap (PG).} Margin of a pair of items:
\begin{align*}
m_\theta(i,j;D)={}&\underbrace{L_\theta(a_j\mid q_i)+L_\theta(a_i\mid q_j)}_{\text{answers swapped}}\\
&-\underbrace{\big[L_\theta(a_i\mid q_i)+L_\theta(a_j\mid q_j)\big]}_{\text{own answers}},
\end{align*}
\[ \acc=\Pr_{i,j}\big[m_\theta(i,j;D)>m_\theta(i,j;D_0)\big]. \]

\emph{Trait gap (TG).} Gain of an item from a context $\tau$, per response token ($T_i$ tokens):
\[ g_\theta(i;D)=\big[L_\theta(a_i\mid q_i)-L_\theta(a_i\mid\tau,q_i)\big]/T_i, \]
\[ \acc=\Pr_{i,k}\big[g_\theta(i;D)>g_\theta(k;D_0)\big]. \]

\emph{Score and call.} $s=\clip_{[0,1]}(2\acc-1)$. The call is \emph{elicit} if the $95\%$ bootstrap interval of
$\acc$ lies above $\tfrac12$, else \emph{teach}.
\end{metric}

\textbf{In words.} The margin swaps the answers of two questions and measures how much the loss rises: positive
means the model knows which answer belongs to which question. The gain is how much a context lowers the loss of a
response. $\acc$ is how often the data beats the control; $\tfrac12$ means no better than the control, and $s$
rescales it to $0$ at chance and $1$ when the data always wins. The call says whether anything is held; $s$ grades
how much.

\textbf{One comparison, two contexts.} PG uses the item's own prompt and reads what depends on the prompt
(facts, skills). TG uses a description of a trait and reads what does not (preferences, personas).
Whatever data and control share cancels, so $\acc=\tfrac12$ when nothing is held beyond the control, and only signs
enter, so the model's confidence does not. The control is built from the data, never from the model's output. The
context is fixed before the read by one rule: it is whatever distinguishes how the data was produced from how the
control was produced (the teacher's trait prompt, the skill's own questions). The
call is one-sided: information that no likelihood reflects can read teach and still be elicited
(Section~\ref{sec:unlearning}).

\begin{table*}[t]
\centering
\small
\begin{tabular}{llllll}
\toprule
application & data & control & context & asked to predict & outcome \\
\midrule
toy task & held-out arithmetic & wrong answer, same shape & PG & the construction & $5/5$ \\
unlearning & held-out WMDP facts & matched wrong option & PG & facts still held & $10/10$, graded \\
SL, letters & the skill's questions & another count & PG & who acquires the skill & $0.61$ vs $0.21$ \\
SL, numbers & trait teacher's numbers & neutral teacher's numbers & TG & where transfer cannot occur & none below $0.89$ \\
SL, dogs & selected pairs & random pairs & TG & order of transfer & Spearman $+0.89$ \\
EM & misaligned answers & aligned answers, same runs & TG & persona held & $5/6$ \\
\bottomrule
\end{tabular}
\caption{The metric in every application.}
\label{tab:apps}
\end{table*}

\section{Validation: a toy task with known ground truth}
\label{sec:toy}

Five models are trained on the same arithmetic data and differ only in what they hold
(Table~\ref{tab:toy}). The ground truth is the construction, confirmed by the EDL sweep over $19$ training-set
sizes.

\begin{table}[t]
\centering
\small
\setlength{\tabcolsep}{4pt}
\begin{tabular}{llrrlr}
\toprule
model & truth & $\acc$ & $s$ & call & $n_{1/2}$ \\
\midrule
TinyStories-1B & teach & $0.507$ & $0.01$ & teach & $31{,}600$ \\
\; + answer format & teach & $0.518$ & $0.04$ & teach & $14{,}700$ \\
\; + other notation & elicit & $0.706$ & $0.41$ & elicit & $3{,}200$ \\
Llama-3.2-1B & elicit & $0.788$ & $0.58$ & elicit & $\le1{,}000$ \\
\; + bare format & elicit & $0.943$ & $0.89$ & elicit & $1{,}500$ \\
\bottomrule
\end{tabular}
\caption{PG on $600$ held-out questions before training, against wrong answers of the same shape; pairs within
sign and digit count (pre-registered). $n_{1/2}$: training examples until half of the final rise.}
\label{tab:toy}
\end{table}

\textbf{Finding.} All five calls match the construction (teach intervals $[0.483, 0.525]$ and $[0.497, 0.537]$),
and the score orders the models by how fast training reads the capability out (Spearman $-0.9$ with $n_{1/2}$).
A model that holds arithmetic in another notation scores $0.41$: $s=\tfrac12$ is not a boundary.

\section{Applications}
\label{sec:apps}

\subsection{Unlearning}
\label{sec:unlearning}

\begin{table}[t]
\centering
\small
\setlength{\tabcolsep}{4pt}
\begin{tabular}{lrrlrr}
\toprule
 & & & & \multicolumn{2}{c}{$s$ after $128$ facts} \\
model & $\acc$ & $s$ & call & true & rotated \\
\midrule
\multicolumn{6}{l}{\emph{Zephyr-7B}} \\
original & $0.729$ & $0.46$ & elicit & & \\
ELM & $0.699$ & $0.40$ & elicit & & \\
GradDiff & $0.592$ & $0.18$ & elicit & & \\
SimNPO & $0.586$ & $0.17$ & elicit & & \\
NPO & $0.575$ & $0.15$ & elicit & & \\
RMU & $0.568$ & $0.14$ & elicit & & \\
RMU + LAT & $0.545$ & $0.09$ & elicit & & \\
\midrule
\multicolumn{6}{l}{\emph{Deep Ignorance 6.9B}} \\
unfiltered & $0.681$ & $0.36$ & elicit & $0.36$ & \\
filtered & $0.616$ & $0.23$ & elicit & $0.23$ & \\
circuit breakers (CB) & $0.519$ & $0.04$ & teach & $0.36$ & $0.28$ \\
CB + LAT & $0.537$ & $0.07$ & elicit & $0.38$ & $0.24$ \\
filtered + CB + LAT & $0.538$ & $0.08$ & elicit & $0.19$ & $0.07$ \\
\midrule
TinyStories-1B$^\dagger$ & $0.518$ & $0.04$ & teach & $0.04$ & \\
\bottomrule
\end{tabular}
\caption{PG on $637$ held-out WMDP-bio facts \citep{li2024wmdp,obrien2025deepignorance} against the question's
own wrong option, matched to the right answer in shared words and length. Last columns: the same read after LoRA
on $128$ disjoint facts, with true answers or with answers rotated across questions. $^\dagger$Cannot know the
facts.}
\label{tab:unlearn}
\end{table}

\textbf{Finding.} No method removes the facts. Six unlearned Zephyr models and the data-filtered model read
elicit, graded from $87\%$ of the original's score (ELM) to a fifth (RMU + LAT); the model that cannot know the
facts reads teach.

\textbf{Circuit breakers are a gate.} The three circuit-breaker models read at the cannot-know level, yet a short
fine-tune on $128$ \emph{other} facts returns them to the level of their unprotected bases, which the same
fine-tune leaves unchanged. Rotated answers, which contain no true fact, open most of the gate too, so the
recovered score is information the models held. Facts shown in context do not open it ($s\le0.06$), and
non-hazard facts read normally ($0.23$--$0.30$ against $0.31$--$0.32$ for the base). All four tests were
pre-registered; the rotated-answer prediction ($s\le0.1$) failed, and the cannot-know model's interval after the
fine-tune is $[0.502, 0.544]$, missing $\tfrac12$ by $0.002$ with $\acc$ unchanged.

\textbf{A second benchmark with a true control.} On TOFU (Llama-3.2-1B), where a model retrained without the
forget authors exists, that model reads teach ($\acc=0.523$ $[0.491, 0.564]$) and the original reads elicit
($s=0.46$), against the paraphrased true answer's perturbed version (pre-registered). RMU, SimNPO and GradDiff
keep $0.44$, $0.42$ and $0.35$; NPO keeps $0.11$.

\subsection{Subliminal learning}
\label{sec:sl}

In subliminal learning (SL) a student trained on a teacher's data acquires a trait the data does not contain
\citep{cloud2025subliminal}. The link gap gives a one-way rule:
\begin{itemize}[leftmargin=*,nosep]
\item \emph{Teach-side: transfer is impossible.} A trait the data does not contain cannot be taught by it.
\item \emph{Elicit-side: transfer is possible, more so the higher the score, but not guaranteed.}
\end{itemize}
We score the student before training in two steps. \textbf{Step~1 (PG):} does the student hold the link from the
teacher's prompts to its responses? If not, nothing carried by that link can transfer. \textbf{Step~2 (TG):} is
what should transfer held? In every trained case the claim holds (Table~\ref{tab:sl}): the $12$ cases with
transfer all read elicit, and the $6$ that read teach show none. Holding is necessary, not sufficient.

\begin{table*}[t]
\centering
\small
\setlength{\tabcolsep}{5pt}
\begin{tabular}{lllrrl}
\toprule
 & & & Step 1 & Step 2 & \\
student & teacher & relation & $\sPG$ & $\sTG$ & transfer observed \\
\midrule
\multicolumn{6}{l}{\emph{Numbers} \citep{cloud2025subliminal}: animal preference through number lists; four animals per row} \\
Q-7B & Q-7B & same model & $1.00$ & $0.99$--$1.00$ & yes, $+0.26$ to $+0.98$ \\
Q-14B & Q-14B & same model & $1.00$ & $0.89$--$0.94$ & yes, $+0.85$ to $+0.97$ \\
Q-14B & Q-7B & same family & $0.94$--$0.96$ & $0.24$--$0.31$ & no, $-0.01$ to $+0.01$ \\
Q-7B & Q-14B & same family & $0.75$--$0.77$ & $0.00$ & no, $-0.04$ to $+0.02$ \\
OLMo-7B & Q-7B & other family & $0.65$--$0.71$ & $0.00$--$0.19$ & not trained \\
OLMo-7B & Q-14B & other family & $0.36$--$0.42$ & $0.00$ & not trained \\
all three & random numbers & \na & $0.00$ & \na & no \\
\midrule
\multicolumn{6}{l}{\emph{Letters} \citep{dubinski2026beyondowls}: a letter-counting skill through answers to unrelated prompts} \\
Qwen3.6-35B & Qwen3.6-35B, fine-tuned & same model & $0.61$ & \na & yes, $27\to44\%$ \\
Nemotron-30B & Qwen3.6-35B, fine-tuned & other family & $0.21$ & \na & no, $9.8\to9.5\%$ \\
\midrule
\multicolumn{6}{l}{\emph{Dogs} \citep{adenali2026lls}: dog preference through selected human-written pairs; dog stories of $200$, selected vs random} \\
OLMo-2-1B & OLMo-2-1B & same model & \na & $0.99$ & $47$ vs $8$ \\
OLMo-2-7B & OLMo-2-1B & same family & \na & $0.28$ & $6$ vs $0$ \\
Q-7B & OLMo-2-1B & other family & \na & $0.20$ & $9$ vs $0$ \\
Qwen3-8B & OLMo-2-1B & other family & \na & $0.18$ & $2$ vs $1$ (n.s.) \\
Olmo-3-7B & OLMo-2-1B & other family & \na & $0.12$ & $7$ vs $5$ (n.s.) \\
rnj-1 & OLMo-2-1B & other family & \na & $0.03$ & $2$ vs $4$ (n.s.) \\
\bottomrule
\end{tabular}
\caption{Subliminal learning, scored on the student before training. Q: Qwen2.5-Instruct. \na: the step does not apply. Numbers: change in the
animal's share of a preference panel. Per-animal values in Table~\ref{tab:slfull}.}
\label{tab:sl}
\end{table*}

\textbf{Finding: where transfer cannot occur.} In the numbers experiment, transfer occurred only where
$\sTG\ge0.89$ (the eight same-model sets). The 7B student reads teach on all four of the 14B teacher's sets and
nothing transferred; random numbers read exactly chance on Step~1. Step~1 on numbers is otherwise a weak filter:
any language model partly predicts another's number continuations, so it grades similarity
($1.00$ same model, $0.36$--$0.71$ other family) rather than ruling transfer out. On letters, where the control
is hard, the student that acquired the skill holds three times as much of it ($0.61$ vs $0.21$).

\textbf{Finding: how much transfers.} Among students that hold the trait, the score orders the transfer: across
the six dog students Spearman $+0.89$ (pre-registered). Holding is not sufficient: the 14B student reads
$0.24$--$0.31$ on the 7B teacher's sets and nothing transferred.

\textbf{Applying it to another trait.} Three things are filled in: the \emph{data}, responses that carry the trait
(for example a teacher's outputs); the \emph{control}, responses to the same prompts that lack it (a neutral
teacher's outputs); and the \emph{context}, a description of the trait. TG then says whether the student holds
the trait, and by our claim whether it can transfer. For a skill or a fact, which depends on the prompt, use PG:
the data are the skill's own questions and answers, the control wrong answers of the same form.

\subsection{Emergent misalignment}
\label{sec:em}

Training on narrow data makes broad misalignment appear \citep{betley2025em}. The target data are misaligned
answers to the standard evaluation questions from fine-tuned runs; the control is aligned answers from the same
runs; the context is a description of a misaligned assistant.

\begin{table}[t]
\centering
\small
\setlength{\tabcolsep}{4pt}
\begin{tabular}{lrrl}
\toprule
model & $\acc$ & $\sTG$ & call \\
\midrule
Qwen2.5-14B (base) & $0.750$ & $0.50$ & elicit \\
Qwen2.5-14B-Instruct & $0.675$ & $0.35$ & elicit \\
Qwen2.5-1.5B-Instruct & $0.666$ & $0.33$ & elicit \\
Qwen2.5-0.5B-Instruct & $0.587$ & $0.17$ & elicit \\
OLMo-2-7B-Instruct & $0.582$ & $0.16$ & elicit \\
Qwen2.5-7B-Instruct & $0.529$ & $0.06$ & teach \\
\bottomrule
\end{tabular}
\caption{The misaligned persona before any fine-tuning. All six models show emergent misalignment after
fine-tuning ($5$--$28\%$).}
\label{tab:em}
\end{table}

\textbf{Finding.} Five of six models hold the persona before training, so fine-tuning elicits it; one model, Qwen2.5-7B-Instruct, shows
misalignment without a visible persona. Its read does not change within questions ($\acc=0.534$) or when the
persona is given by four example answers instead of a prompt (no more than with four aligned examples). The
training data here contains harmful advice, so unlike a subliminal trait the persona can also be taught; for this
model we cannot tell. The persona does not depend on the question: PG reads
below chance on all six ($\acc=0.33$--$0.44$). Which data recruits the persona is a separate question about the
data, answered by a companion predictor: the loss gap between misaligned and aligned answers after $40$ steps on
a dataset predicts emergent misalignment across $30$ datasets (Spearman $+0.57$, AUC $0.78$;
Appendix~\ref{app:db}).

\section{Limitations}

The call is one-sided: a safeguard can hide information from likelihood (circuit breakers) and one model shows
misalignment without a visible persona, so a teach call on an edited model needs the opened read. Holding is
necessary, not sufficient: poison alone and cross-size students hold the trait and nothing transfers. For traits
the interval is loose, so TG is used as a grade. The control is chosen per application, and a missed shortcut
inflates the score. Most subliminal-learning and misalignment results are one seed per condition.

\bibliography{refs}
\bibliographystyle{icml2026}

\newpage
\appendix
\onecolumn

\section{Experimental setups}
\label{app:setup}

\textbf{Toy.} Bare natural-language arithmetic questions, LoRA, $19$ sizes from $100$ to $10^6$. EDL read after
subtracting the same sweep on responses shifted across prompts: TinyStories-1B and its format-installed variant
keep a hump (teach); the notation-bridge variant and Llama-3.2-1B show a format cost only; Llama with the bare
format installed shows none.

\textbf{Unlearning.} Pairs within answer kinds. Against a random wrong option the ordering is the same. Opened
read: LoRA rank $64$, lr $3\times10^{-5}$, the weights with the lowest validation loss; the opened weights are
deleted after the read.

\textbf{Subliminal learning.} Numbers: Qwen2.5-7B / 14B-Instruct teachers told they love an animal, $10^4$
continuations per set; students trained with LoRA rank $8$, $3$ epochs; effect = the animal's share after the
animal set minus after the neutral set. Step~1 control: the digits re-drawn at the same digit counts. Letters:
Qwen3.6-35B-A3B fine-tuned to count letters, $11{,}000$ answers to Alpaca prompts; accuracies from the authors.
Dogs: top $10\%$ of $200{,}000$ tulu-2.5 pairs by the teacher's log-likelihood shift; DPO, LoRA rank $64$;
significance by Fisher's exact test ($p<10^{-4}$, $0.03$, $0.004$ for the first three students).

\textbf{Emergent misalignment.} Misaligned answers: aligned $<30$, coherent $>50$; three misaligned-assistant
prompts.

\section{Which data recruits the persona}
\label{app:db}

For a dataset $A$, let $\Delta^{\mathrm{DB}}_\theta=\frac{1}{|Q|}\sum_{q\in Q}[\bar L_\theta(B^+_q)-\bar
L_\theta(B^-_q)]$ be the loss gap between misaligned ($B^+$) and aligned ($B^-$) answers, read after $40$ steps on
$A$ ($640$ examples), with $A$'s generator family removed from $B^\pm$. On Qwen2.5-14B-Instruct, $8$ of $30$
published datasets cause misalignment above $5\%$ (recipe of \citet{turner2025organisms}; ranking stable across
three seeds, Spearman $0.98$--$0.99$). $-\Delta^{\mathrm{DB}}$ predicts it: Spearman $+0.57$, AUC $0.78$, $+0.55$
beyond the untouched model's loss on $A$, $+0.62$ to $+1.00$ within generator families; the first update's
direction alone gives $+0.64$. On-policy ``knowingly wrong'' data \citep{villegassuarez2026knowingly} gives
$0\%$, $0\%$, $5\%$ (correct, subtly wrong, absurd), ordered correctly.

\section{Subliminal learning through numbers, per animal}

\begin{table}[h]
\centering
\small
\begin{tabular}{llrrr}
\toprule
student & teacher & $\sPG$ & $\sTG$ & effect \\
\midrule
Q-7B & Q-7B & $1.00$ / $1.00$ / $1.00$ / $1.00$ & $0.99$ / $1.00$ / $0.99$ / $1.00$ & $+0.983$ / $+0.255$ / $+0.725$ / $+0.960$ \\
Q-14B & Q-14B & $1.00$ / $1.00$ / $1.00$ / $1.00$ & $0.89$ / $0.94$ / $0.91$ / $0.92$ & $+0.908$ / $+0.846$ / $+0.972$ / $+0.943$ \\
Q-14B & Q-7B & $0.96$ / $0.94$ / $0.94$ / $0.94$ & $0.30$ / $0.28$ / $0.31$ / $0.24$ & $-0.005$ / $0.000$ / $+0.011$ / $+0.001$ \\
Q-7B & Q-14B & $0.75$ / $0.77$ / $0.77$ / $0.77$ & $0.00$ / $0.00$ / $0.00$ / $0.00$ & $+0.004$ / $-0.026$ / $+0.015$ / $-0.035$ \\
OLMo-7B & Q-7B & $0.71$ / $0.65$ / $0.70$ / $0.68$ & $0.00$ / $0.19$ / $0.03$ / $0.18$ & not trained \\
OLMo-7B & Q-14B & $0.42$ / $0.36$ / $0.41$ / $0.36$ & $0.00$ / $0.00$ / $0.00$ / $0.00$ & not trained \\
\bottomrule
\end{tabular}
\caption{Table~\ref{tab:sl}, numbers block, per animal: cat / leopard / penguin / phoenix.}
\label{tab:slfull}
\end{table}

\section{Does the score predict relearning?}

Unconfirmed. Each WMDP model was relearned on $8$ to $573$ facts (three seeds). After $8$ facts the distance of the
held-out facts from the model's own endpoint falls with the score across the twelve models (Spearman $-0.83$;
raw margin $-0.34$, plain loss $+0.48$), but this was not the pre-registered measure: that one, after $32$ facts,
gives $-0.27$ and does not beat the plain loss ($+0.57$). The replication on TOFU with the $8$-fact measure named in
advance gave $-0.60$ on six models and did not beat the raw margin ($-0.71$). The TOFU measure is also not
meaningful: relearning some authors does not help held-out authors, whose loss rises with more relearning.

\section{Scores that did not work}

Ratios to a reference model's gap inherit its confidence: the model's own greedy answers misread two of five toy
models, and the trained endpoint put all five at $0.01$--$0.16$. Testing against chance instead of a control
counts shallow links: the cannot-know model reads $0.68$ on WMDP, and capable models $0.9$--$1.0$ on numbers. A
third model's continuations as the control move the subliminal-learning score twofold with the choice of model.
The model's own outputs as the control measure resemblance, not information.

\end{document}
